%%----------------------------------------------------------------------------%%
%% Student-specific compatibility macros
%%
%% This file is only for packages and macros needed by this student's chapters.
%%
%% Do not load hyperref, caption, titlesec, trackchanges, tcolorbox, or
%% \begin{document} here. Those belong in the main template files.
%%
%% Theorem-like environments are defined in ncsu-required.tex. Optional visual
%% changes to those environments belong in optional-accessible.tex.
%%----------------------------------------------------------------------------%%

%%----------------------------------------------------------------------------%%
%% Packages needed by the student's content
%%----------------------------------------------------------------------------%%

\usepackage{accents}
\usepackage{comment}
\usepackage{easytable}
\usepackage{multirow}
\usepackage{multicol}

%% This package lets us catch a few literal Unicode characters that may appear
%% in the student's source. It is especially useful for avoiding PDF/A .notdef
%% glyph failures from unsupported symbols.
\usepackage{newunicodechar}

%%----------------------------------------------------------------------------%%
%% TikZ libraries
%%
%% If you do not use TikZ, delete this block.
%% Keep these centralized here so individual figure files do not have to load
%% TikZ libraries themselves.
%%----------------------------------------------------------------------------%%

\usepackage{tikz}
\usetikzlibrary{
  arrows,
  arrows.meta,
  backgrounds,
  calc,
  decorations.pathmorphing,
  decorations.pathreplacing,
  fit,
  math,
  positioning,
  quotes,
  through
}

%% If the original TikZ code used uppercase color names as TikZ keys, these
%% definitions prevent errors such as: I do not know the key '/tikz/Red'.
%%
%% These are RGB-based color aliases, not CMYK definitions.

\definecolor{AccessibleYellow}{HTML}{FFF4B8}
\definecolor{AccessibleCyan}{HTML}{D8F3F5}
\definecolor{AccessibleRed}{HTML}{F7D6D0}
\definecolor{AccessibleGreen}{HTML}{DCEFD8}
\definecolor{AccessibleGray}{HTML}{F2F2F2}
\definecolor{AccessibleLightGreen}{HTML}{EAF6EA}
\definecolor{AccessibleLightBlue}{HTML}{EAF3FA}
\definecolor{AccessibleMediumBlue}{HTML}{BFDDF2}
\definecolor{AccessibleBlue}{HTML}{D9EAF7}
\definecolor{AccessibleLightCyan}{HTML}{D8F3F5}

%%----------------------------------------------------------------------------%%
%% Basic student macros
%%----------------------------------------------------------------------------%%

\providecommand{\uv}[1]{\ensuremath{\mathbf{\hat{#1}}}}
\providecommand{\bo}{\ensuremath{\mathbf{\Omega}}}

\providecommand{\eref}[1]{Eq.~\ref{#1}}
\providecommand{\fref}[1]{Fig.~\ref{#1}}
\providecommand{\tref}[1]{Table~\ref{#1}}
\providecommand{\cref}[1]{Chap.~\ref{#1}}

\providecommand{\todo}[1]{{\color{red}\textbf{TODO: #1}}}

\providecommand{\R}{\mathbb{R}}

%%----------------------------------------------------------------------------%%
%% Other student-specific additions
%%----------------------------------------------------------------------------%%
